\documentclass[11pt]{article}
\usepackage[T1]{fontenc}
\usepackage{lmodern}
\usepackage[margin=1in]{geometry}
\usepackage{amsmath,amsthm,microtype}
\usepackage[hidelinks]{hyperref}

\newtheorem{theorem}{Theorem}
\hypersetup{
  pdftitle={Small maximal arrangements in Euler's 36 officers problem},
  pdfauthor={Aafi Khan Malek}
}

\begin{document}
\begin{center}
{\Large\bfseries Small maximal arrangements in Euler's 36 officers problem\par}
\medskip
Aafi Khan Malek\\
Independent researcher\\
\href{mailto:aafimalek2032@gmail.com}{\texttt{aafimalek2032@gmail.com}}\par
\smallskip
26 November 2026
\end{center}

\begin{abstract}
Euler's classical problem asks for an arrangement of all 36 officers
and has a negative answer. We consider the fewest officers needed
to make a partial arrangement maximal. We prove that the minimum is
18 among arrangements whose graph of used pairs is disconnected.
Any smaller maximal arrangement must therefore be connected. We also exclude all maximal arrangements of 14 or 15 officers, so the unrestricted minimum lies between 16 and 18. These results use independently checked SAT certificates.
\end{abstract}

\section{Problem and result}

There are six ranks and six regiments, one officer of each rank
and regiment, and a \(6\times 6\) grid. Each row and each column
must use distinct ranks and distinct regiments. A placement is
\emph{maximal} if no unused officer can be placed in an empty
cell under these rules. In the language of orthogonal partial
Latin squares, the minimum was left unresolved at order~6
by Donovan, Grannell and Yaz\i c\i~\cite[Table~1]{dgy}.
Euler's question about placing all 36 officers has a negative
answer~\cite{euler}.

Represent each placement by a word specifying the row, column,
rank and regiment. Distinct words agree in at most one coordinate.
Construct a graph on the 24 labels, six per coordinate, with
an edge between two labels if a word contains both. We call it
the used-pair graph. Maximality means that every word, whether
already placed or absent, contains an edge of this graph.

\begin{theorem}
If the used-pair graph of a maximal partial arrangement is
disconnected, the arrangement has at least 18 officers. This
bound is attained.
\end{theorem}

\section{Reduction to two orders}

Suppose there are at most 17 words. Every label occurs. Indeed,
if a label is unused, the projection onto the other three
coordinates must be a maximal partial Latin square of order 6.
Otherwise an addable triple together with that unused label gives
an addable word. The projection has the same number of entries,
since two words cannot agree in three coordinates. The known
minimum for a maximal partial Latin square of order 6 is
\(6^2/2=18\)~\cite[Section~1]{dgy2}.

Every component of the used-pair graph consequently meets all
four coordinates. To see this, every vertex lies in a word,
and that word induces a clique containing one label from
each coordinate. There are at most three components. Four would
let us choose one label in each coordinate from distinct
components, producing an addable word.

If there are three components, then every pair within a component
from distinct coordinates must be used. An unused pair could be
joined by a label from each of the two other components to make an
addable word. Completeness between the coordinate groups of one
component forces the four group sizes to be equal, since each
cross-coordinate projection pairs all labels and has the same
number of words. Thus the three components are full orthogonal
pairs of orders summing to 6. No pair of orthogonal Latin squares
has order 2, so the only possible partition into three positive
orders is \(1,1,4\). It uses \(1+1+16=18\) words.

Now suppose there are two components \(A,B\). Let \(a_i\) and
\(b_i=6-a_i\) be their sizes in coordinate \(i\), and let \(U,V\)
be their numbers of words. For any pair of coordinates \(i,j\)
we have \(U\le a_i a_j\) and \(V\le b_i b_j\). Reorder the
coordinates so \(a_1\le a_2\le a_3\le a_4\). If the sizes vary,
then \(a_1a_2<a_3a_4\) and \(b_3b_4<b_1b_2\). There would
be an unused pair in \(A\) on coordinates \(3,4\), and another
in \(B\) on \(1,2\). Combining them violates maximality.
Thus \(a_i=a\) and \(b_i=b\) for all \(i\).
At least one component is full, since otherwise an unused pair
in each component on complementary coordinate pairs again
gives an addable word. Say \(B\) is full of order \(b\).

Each projection of \(A\) onto three coordinates is a maximal
partial Latin square of order \(a=6-b\). Otherwise an unused
triple in \(A\) could be extended by a label in \(B\).
The value \(b=2\) is impossible for a full orthogonal pair.
The value \(b=5\) already uses 25 words, and \(b=4\) uses 16
with at least 2 more in \(A\), by the order-2 bound for
maximal partial Latin squares. The remaining cases are
\[
(a,b)=(3,3),\ U\le8
\qquad\text{or}\qquad
(a,b)=(5,1),\ U\le16.
\]
We exclude both by the finite checks below. Conversely,
two disjoint full orthogonal pairs of order 3 give a
disconnected maximal arrangement with 18 words.
Every possible word repeats a block in two coordinates,
and that pair is used in the block.

\section{Finite exclusions}

For a candidate design of order \(a\) whose every three-coordinate
projection is maximal, let \(m\) be its minimum label frequency.
Fix a label of frequency \(m\) in one of the three coordinates.
In each other coordinate it uses \(m\) labels. All pairs in the
remaining \((a-m)\times(a-m)\) rectangle must be used, by maximality
of the projection. The sum of frequencies in one coordinate
is at least \((a-m)^2+m^2\). Also \(m\le\lfloor U/a\rfloor\).
These bounds leave \(m\in\{1,2\}\) for \(a=3,U\le8\) and
\(m\in\{2,3\}\) for \(a=5,U\le16\).

The accompanying generator makes one formula for each case.
It uses a variable for every possible word and every
cross-coordinate pair. Each used pair has exactly one supporting
word. The sum of used pairs in any fixed projection is the word
count. Every three-coordinate tuple must contain a used pair.
Unary frequency thresholds are defined by the exact truth table of
the used neighbors in one other coordinate.

The generator fixes a minimum-frequency label as label 0 of the
first coordinate. Permute the other labels independently to align
its \(m\) incident words as \((0,i,i,i)\) for \(0\le i<m\).
Ordering the remaining labels by frequency in each coordinate
respects those fixed words. The aligned labels can be permuted
jointly to order their frequencies in the second coordinate.
For an unused pair with frequency bounds \(d,e\) on its endpoints,
each of the two remaining coordinates has at least \(a-d-e\)
candidate values. All pairs between the candidate values must be
used. Writing \(r=a-d-e>0\), the used pairs force at least
\(r^2+(a-r)m\) words by counting in one coordinate. The generator
uses this necessary condition to prune clauses.

\begin{center}
\begin{tabular}{rrrrl}
\(a\) & Word limit & \(m\) & Variables & Certificate\\
\hline
3 & 8 & 1 & 615 & Verified\\
3 & 8 & 2 & 615 & Verified\\
5 & 16 & 2 & 4639 & Verified\\
5 & 16 & 3 & 4639 & Verified
\end{tabular}
\end{center}

CaDiCaL produced the proofs. DRAT-trim checked the proofs and
extracted smaller traces. The public archive includes those traces
and the generator to check them against the exact formulas.
The \(18\)-officer bound is exact for disconnected arrangements.
The unrestricted minimum is still unresolved by this argument.

\section{Excluding fourteen officers}

\begin{theorem}
Every maximal partial arrangement has at least \(15\) officers.
\end{theorem}

Write \(F\) for the number of officers in a maximal arrangement and
\(m\) for the minimum label frequency. The frequency bound of
Donovan, Grannell and Yaz\i c\i~\cite[Eq.~(1)]{dgy2} gives \(F\ge12\).
At \(F=12\), the equality structure requires three full pairs of
order 2~\cite{gap}; these do not exist. The same note excludes
\(F=13\). We assume \(F=14\) and obtain another contradiction.

\paragraph{Minimum frequency one.} By the same frequency bound and
the average frequency \(14/6\), \(m\in\{1,2\}\). If \(m=1\),
fix a frequency-one label \(x\). For any two other coordinates
take a maximum matching of unused pairs among the five labels in
each coordinate unused with \(x\). If the matching size is \(t\),
the bound gives
\[
28\ge 2+25+(5-t)^2.
\]
Thus \(t\ge4\). If \(t=5\), every one of the five labels unused
with \(x\) in the fourth coordinate must meet every edge of the
matching through a used pair at one endpoint. Each such label
has frequency at least three, since it can meet at most two
matching edges per occurrence. Counting the fourth coordinate
gives \(F\ge 5\cdot3+1=16\), a contradiction. Therefore \(t=4\)
and the matching bound is tight. Its equality condition forces
the two neighbors of \(x\) in the two matched coordinates to
have frequency one. The matched coordinates are arbitrary, so
all three neighbors of \(x\) have frequency one. The unique word
containing \(x\) is then a component by itself. This contradicts
the disconnected bound just proved.

\paragraph{Minimum frequency two.} Call a frequency-two label
low and any other label high. Each coordinate has at most two
high labels because its total excess over frequency two is two.
There are two cases.

First suppose an unused pair of low labels in distinct coordinates
has a common used neighbor in another coordinate. Put the two
labels and the shared neighbor at value 0 in coordinates 1, 2,
and 3, respectively. The unused pair leaves at least three
candidate labels in coordinate 3 and at least two in coordinate 4.
Every pair between the candidate sets is used. Consequently
the fourth-coordinate candidate set consists of exactly two
high labels of frequency three, and the third-coordinate candidate
set has exactly three labels. The two low labels' neighbor sets
in coordinate 3 meet only at 0. In coordinate 4 the two sets
are disjoint and contain the four low labels.

Permute the labels in the other positions to align the four
incident words, with zero-based labels, as
\[
(0,1,0,0),\quad(0,2,1,1),\quad(1,0,0,2),\quad(2,0,2,3).
\]
This leaves labels 4 and 5 high in coordinate 4 and labels 3, 4,
and 5 as candidates in coordinate 3. All six candidate pairs
are required. The SAT formula imposes these facts. In each
coordinate it also enforces the degree bound four and total
excess two. Sorting the frequencies of the values 3 through 5
in the first three coordinates is allowed by permutations that
preserve the four fixed words and the required pairs.

Now suppose that every pair of low labels in distinct coordinates
with a shared used neighbor is itself used. If every label has
a low neighbor in every other coordinate, the low-component
argument in the preceding note applies~\cite[Section~3]{gap}.
It makes the graph on the low labels complete across the coordinate
parts of each component. Each component meets every coordinate,
with at most two low labels in each. Each high label is attached
to the unique component containing all its low neighbors.
Every component with the attached labels has at least two
labels in each coordinate, because a low label has both its
neighbors in that component. There are at least two components
because each coordinate has at least four low labels, and at
most three because each component uses at least two labels.

With two components, each part of each component already contains
two low labels, so every low label has its two neighbors in each
other coordinate among the low labels. This leaves the high labels
without low neighbors. With three components, each part has
exactly two labels after attachment. There are at most 12 internal
words, and any word crossing the components has at least one
same-component pair. That used pair further reduces the internal
capacity to 11. Thus at least three words cross the components,
and all of them consist of high labels. Among three such words
each of the four coordinates forces an agreement between two words,
so some pair of words agrees twice. This contradicts the Latin
and orthogonality rules.

There must therefore be a label with no low neighbor in another
coordinate. It has frequency two and its two neighbors in that
coordinate are the high labels, each of frequency three. Normalize
its two incident words to \((0,0,0,0)\) and \((0,1,1,1)\), with
the second-coordinate values 0 and 1 high. The formula in this
case also enforces the assumed common-neighbor property on low
labels. Every candidate 14-officer arrangement with minimum
frequency two can be relabeled into one of the two formulas.

\begin{center}
\begin{tabular}{lrrl}
Case at \(F=14\) & Variables & Clauses & Certificate\\
\hline
Unused low pair with a shared neighbor & 10016 & 40260 & Verified\\
Common-neighbor property and a missing low neighbor & 9856 & 42524 & Verified
\end{tabular}
\end{center}

Both formulas are unsatisfiable by the supplied DRAT certificates.

\section{Excluding fifteen officers}

\begin{theorem}
Every maximal partial arrangement has at least \(16\) officers.
\end{theorem}

Suppose \(F=15\). If a label had frequency zero, the frequency
bound would give \(2F\ge 6^2\). The average frequency is \(15/6\),
so \(m\in\{1,2\}\).

\paragraph{Minimum frequency one.}
Relabel the unique word containing a minimum label as
\((0,0,0,0)\), with the minimum label in coordinate 1. For any two
other coordinates, let \(u,v\) be the frequencies of the two
labels of this word in those coordinates. Keep the two
excluded-frequency sums explicit in the matching argument
of~\cite[Theorem 2.1]{dgy2}. They are \(u\) and \(v\), since
the minimum label has exactly one neighbor in each coordinate.
The bound becomes
\[
2F\ge u+v+25+(5-t)^2.
\]
A perfect matching (\(t=5\)) would need at least 16 words by the
same argument as at \(F=14\). Hence \(t\le4\) and \(u+v\le4\).
This holds for any two of the three neighbors. The word cannot
have all four frequencies equal to one, since it would be an
isolated component. Up to permutation of the other coordinates,
the three neighbor frequencies are therefore
\[
(1,1,2),\quad(1,1,3),\quad(1,2,2),\quad(2,2,2).
\]

\paragraph{Minimum frequency two.}
Every coordinate has three units of excess above frequency two.
Each contains at least one low label. Choose coordinate 1 with the
fewest high labels, say \(p\), where \(p\in\{1,2,3\}\).
Put its high labels last. When \(p=2\) their frequencies are
three and four. In that case, put the unique frequency-four label
at value 5 and align its four words as \((5,a,a,a)\) for
\(a=0,1,2,3\). In the other cases, align the two words through a
frequency-two label instead. Linearity makes the labels of the
aligned words distinct in each other coordinate.

The seven formulas include a saturation clause for every possible
word and enforce the excess sum in each coordinate. As in Section~3,
the free suffix labels are ordered by frequency. Among equally
frequent free labels, the generator orders their used-pair vectors.
Take the lexicographically first list of incidence matrices for
coordinate pairs \((1,2),(1,3),(1,4)\) over the permitted
relabelings, with each matrix read row by row. Swapping two
adjacent, equally frequent free rows of the first matrix, or two
equally frequent free columns of any matrix, would improve this
list if their vectors were out of order. The remaining aligned
labels may be permuted jointly to order their frequencies in
coordinate 2. In the frequency-four case its pinned label is
unique, so the row swaps in coordinate 1 preserve every anchor.

\begin{center}
\begin{tabular}{@{}lrrl@{}}
Case at \(F=15\) & Variables & Clauses & Certificate\\
\hline
\(m=1\), each of four frequency triples & 11382 & 53826 & Verified\\
\(m=2\), \(p=1\) & 10374 & 46367 & Verified\\
\(m=2\), \(p=2\) & 10374 & 44483 & Verified\\
\(m=2\), \(p=3\) & 10374 & 46423 & Verified
\end{tabular}
\end{center}

All seven clause systems are unsatisfiable by the archived DRAT
certificates. Together with the 18-officer construction, this gives
\[
16\le \min F\le18.
\]
Searches of the public literature for this order-6 problem and the
related spectrum question found the published 18-officer
construction and no earlier proof of the 14- or 15-officer
exclusions. Horak and Kim mention work by Hor\'ak, Meszka and Rosa
on the related spectrum, so priority remains open to correction~\cite{room}.


\begin{thebibliography}{9}
\bibitem{dgy} D. M. Donovan, M. J. Grannell and E. \c{S}. Yaz\i c\i,
\emph{On maximal orthogonal partial Latin squares and minimal codes with
specified length, minimum distance and covering radius},
Des. Codes Cryptogr. \textbf{93} (2025), 5099--5113.
\href{https://doi.org/10.1007/s10623-025-01704-x}{doi:10.1007/s10623-025-01704-x}.

\bibitem{dgy2} D. M. Donovan, M. Grannell and E. \c{S}. Yaz\i c\i,
\emph{On the minimum number of entries in a pair of maximal orthogonal
partial Latin squares}, \href{https://arxiv.org/abs/2602.09908}{arXiv:2602.09908} (2026).

\bibitem{euler} D. Betten,
\emph{A direct proof of the non-existence of a MOL(6)},
\href{https://doi.org/10.1007/s00022-023-00669-3}{doi:10.1007/s00022-023-00669-3} (2023).
\bibitem{gap} A. K. Malek,
\emph{A gap above the \(n^2/3\) bound for maximal orthogonal partial Latin squares},
\href{https://github.com/Aafimalek/novel-discoveries/tree/main/orthogonal-latin-square-gap}{research note} (2026).

\bibitem{room} P. Horak and D. Kim,
\emph{An improved lower bound for maximal partial Room squares},
\href{https://arxiv.org/abs/2611.07117}{arXiv:2611.07117} (2026).
\end{thebibliography}

\noindent\small OpenAI Codex assisted with the proof, computation,
literature search, and preparation of this note.
\end{document}
